\section{Arithmetic certificates and formalization}\label{app:formalization}
This appendix specifies the integer model used in the classification and
repair proofs. Its domain and entry rules are uniform in $h$, so the two
certificate lemmas have a single depth parameter.

\subsection{The integer block domain}\label{app:integer}
For an integer $h\ge1$, the admissible records are of three kinds:
\begin{align*}
 M(j,i,g)&=(0,j,i,g,0,0,0),\\
 U(d,p,s)&=(1,0,0,0,d,p,s),\\
 D_*&=(2,0,0,0,0,0,0).
\end{align*}
The seven fields are, in order, kind, mode, marker level, marker group,
variable depth, node, and sign. A marker satisfies
$0\le j<6$, $1\le i\le h$ and $g\in\{1,2\}$. The mode codes are
\[
 0=V^+,\quad 1=V^-,\quad 2=W_0^+,\quad 3=W_1^+,
 \quad 4=W_0^-,\quad 5=W_1^-.
\]
A variable satisfies $0\le d\le h$, $p\ge0$, and $d=0\Rightarrow p=0$.
Its sign is $s=2$ when $d=h$, and $s\in\{0,1\}$ when $d<h$, with
$0,1,2$ denoting plus, minus and leaf, respectively.

The physical encoding assigns the appropriate mode code to a marker,
copies the depth, node and sign of a variable, and sends the dummy to
$D_*$. Thus every physical block gives an admissible record. The integer
domain relaxes physical validity by omitting the bound $p<2^d$.

\subsection{Comparison and entry predicates}
For markers, comparison is lexicographic on
\[
 \bigl(\ind[j<4],i,j,g\bigr).
\]
This places modes 4 and 5 first, in precisely the marker order of
Section~\ref{sec:host}. Every marker precedes every non-marker. Among
non-markers the dummy follows variables on rows and precedes them on
columns.

For variable records, write $(o_r,c_r)=(1,0)$ and
$(o_c,c_c)=(0,1)$ for the opening and closing signs on the two axes.
Put $q'=\lfloor q/2\rfloor$ and $p'=\lfloor p/2\rfloor$. The integer
comparison $\prec_a$, for $a\in\{r,c\}$, is defined by
\begin{equation}\label{eq:integer-order}
 U(d,p,s)\prec_a U(e,q,t)\ \Longleftrightarrow\
 \begin{cases}
 p<q\ \lor\ (p=q\land s=o_a\land t=c_a),&d=e,\\
 p<q'\ \lor\ (p=q'\land s=o_a),&e=d+1,\\
 p'<q\ \lor\ (p'=q\land t=c_a),&d=e+1,\\
 \mathrm{true},&|d-e|\ge2.
 \end{cases}
\end{equation}
At equal depths and adjacent depths these are the exact comparisons of
the tree blocks. At larger depth gaps the relation permits every pair.

Let $E_h^0$ be the original integer entry function. Its marker--marker
rule is
\[
 E_h^0(M(j,i,g),M(k,\ell,u))=\ind[j=k\land i=\ell\land g\ne u].
\]
The marker roles are those of the table in Section~\ref{sec:host}, with
level and parity tested on the records. In a signed role, $U(d,p,2)$
belongs to either sign when $d=h$. Entries with a dummy and a non-marker
are one. The variable entries are exactly \eqref{eq:values}, with the
depth and sign conventions above. Let $E_h^1$ have the same rules, except
\[
 E_h^1(D_*,U(h,p,2))=0\qquad(p\ge0).
\]
These definitions retain every entry rule while relaxing node ranges and
some comparisons.

For $\eta\in\{0,1\}$, define $\mathcal C_h^\eta(r,c)$ for two
four-tuples of records by the following conditions: all eight records are
admissible, $r_i\prec_r r_j$ and $c_i\prec_c c_j$ for every
$1\le i<j\le4$, and
\[
 E_h^\eta(r_i,c_j)=(H_4)_{ij}\qquad(1\le i,j\le4).
\]
In particular, all six comparisons on each axis are part of the copy
predicate.

\begin{lemma}[Encoding a physical copy]\label{lem:encoding}
The physical block encoding maps every ordered block copy in $A_h$ to
a witness for $\mathcal C_h^0$, and every ordered block copy in $B_h$ to
a witness for $\mathcal C_h^1$.
\end{lemma}
\begin{proof}
Physical marker levels and groups satisfy the integer ranges. For a
variable, physical validity gives its depth and sign, the nonnegative
node label, and the root condition. Thus the encoding preserves validity.

For comparisons, a marker pair has the same lexicographic key on both
axes. A marker--non-marker pair and a dummy--variable pair follow the
stated axis rules. For two variables at equal depth, node order and the
opening and closing signs give the first clause of
\eqref{eq:integer-order}. At adjacent depths the ancestor node is the
integer quotient by two, giving the second and third clauses. At a gap
of at least two levels the relaxed clause is satisfied. Hence all physical
comparisons imply their integer counterparts, separately for rows and
columns.

Finally, marker pairs use the same mode and group test, mixed marker
pairs use the same role test, and the non-marker entries use the same
dummy and variable rules. The canonical change is precisely on the
encoded dummy-row--leaf-column products. All sixteen entry equalities
are therefore preserved. These arguments are the encoding lemmas
\code{encode_valid}, \code{encode_before},
\code{encode_original_entry} and \code{encode_repaired_entry}; together
they give \code{encode_copy}.
\end{proof}

\subsection{The certificate statements}
Define $T_h(r,c)$ to mean that there exist $j\in\{4,5\}$ and $p\ge0$
such that
\begin{align*}
 r&=\bigl(M(j,h,1),M(j,h,2),U(h,p,2),D_*\bigr),\\
 c&=\bigl(M(j,h,1),v,U(h,p,2),U(h-1,\lfloor p/2\rfloor,1)\bigr),
\end{align*}
where either $v=M(j,h,2)$, or $v=U(h,q,2)$ with
$0\le q<p$ and $q\equiv p\pmod2$. The marker parity is recovered
from the entry conditions when this predicate is decoded in the physical
host.

\begin{lemma}[Arithmetic classification and repair]\label{lem:arithmetic}
For every integer $h\ge1$ and every pair of four-tuples $r,c$ of integer
records,
\begin{equation}\label{eq:certificate-claims}
 \mathcal C_h^0(r,c)\ \Longrightarrow\ T_h(r,c),
 \qquad
 \mathcal C_h^1(r,c)\ \Longrightarrow\ \bot.
\end{equation}
\end{lemma}
\begin{proof}[Certificate proof]
The first implication is the Lean theorem
\code{Arithmetic.classification}, and the second is
\code{Arithmetic.repaired_no_copy}. The classification derivation has
266 source segments and the repair derivation has 243. Their steps are
propositional inferences, linear integer arithmetic lemmas, and
congruence proofs. The transparent formula abbreviations express the
admissibility, comparison and entry predicates specified above.

The two bridge modules instantiate the final derivations with the fields
of $r$ and $c$. Both copy predicates supply eight validity conditions,
twelve pairwise comparisons and all sixteen entry conditions. The
classification derivation needs only fifteen of the entry conditions:
it does not use the condition at $(0,0)$ (with zero-based indices).
The repair derivation uses all sixteen. For classification the negation
of $T_h$ gives the counterexample premise; for repair the premises are
the repaired copy conditions themselves. Thus the omitted classification
premise strengthens that certificate without weakening the definition
of an induced copy.
The resulting contradictions establish \eqref{eq:certificate-claims}.
The derivations and the bridge proofs are checked by the Lean kernel.
\end{proof}

To decode the first conclusion, the leaf-row entry against the first
marker column is one. The role test consequently gives
$j=4+(p\bmod2)$. Physical validity then supplies $p<2^h$ and the valid
minus parent. This is the theorem \code{terminal_decodes}, which yields
the terminal families in Section~\ref{sec:classification}.

\subsection{Formal statements and verification}
The matrix family \code{standardHost} uses the usual order on
$\operatorname{Fin}(n_h)$ on both axes. Its copy predicate
uses increasing four-tuples and all sixteen Boolean equalities. Its
edit distance is the attained minimum Hamming cost over pattern-free
binary matrices. The principal mathematical interfaces are
\begin{center}\small
\begin{tabular}{p{47mm}p{99mm}}\toprule
Result&Lean declaration\\\midrule
Physical size and order&\code{physicalPairs_length}, \code{orderedPosition_order_iff}\\
Matrix copy bridge&\code{standard_copy_iff}\\
Terminal classification&\code{classification_iff}\\
Repaired freeness&\code{canonical_standard_free}\\
Exact copy count&\code{exact_copy_count}\\
Arbitrary-edit lower bound&\code{arbitrary_edit_lower_bound}\\
Exact edit distance&\code{exact_edit_distance}, \code{exact_normalized_distance}\\
Asymptotic ratio&\code{normalized_copy_ratio_tendsto}\\
Removal obstruction&\code{no_polynomial_removal_bound}\\
Main theorem&\code{end_to_end_main_theorem}\\\bottomrule
\end{tabular}
\end{center}

The XOR construction is developed in \code{XorSampling.lean}; its counting
argument is in \code{FiniteAvoidance.lean}. Terminal parameters are counted
in \code{TerminalCount.lean}, and their physical realizations are
enumerated in \code{CopyFibres.lean}. The module \code{GlobalCounting.lean}
identifies their disjoint union with the full copy enumeration.
\code{EditDistance.lean} transports Hamming cost through the
order-preserving bijections and proves attainment of the minimum.

The project pins Lean 4.24.0 and Mathlib at revision
\begin{center}\code{f897ebcf72cd16f89ab4577d0c826cd14afaafc7}.\end{center}
The clean verification procedure is run from the project root by
\begin{quote}\ttfamily\small
bash verify\_all.sh
\end{quote}
Its declaration audit selects every definition and theorem by its defining
project module, including private and anonymous declarations, and checks
the complete axiom dependency closure. Only \code{propext},
\code{Classical.choice} and \code{Quot.sound} are allowed. Current run
records, exact source and artifact hashes, and dependency provenance are
under \code{evidence/strict_build/}; supplied older logs are preserved
separately as historical evidence. Appendix~\ref{app:reproduce} describes
the build procedure, proof dependencies and retained verification records.
